%
% Halaman Abstract
%
% @author  Andreas Febrian
% @version 1.00
%

\chapter*{\normalsize ABSTRACT}
\begin{singlespace}
\vspace*{0.2cm}

\noindent \begin{tabular}{l l >{\raggedright\arraybackslash}p{11.0cm}}
	Name&: & \penulis \\ % ------------ Jangan diedit
	Program&: & \jenjanginggris~Study Program of~\bidanginggris \\ % ------------ Jangan diedit
	Title&: & \judulInggris \\ % ------------ Jangan diedit
\end{tabular}

\vspace*{0.5cm}

\noindent

The development dynamics in the strategic regions of Banten, DKI Jakarta, and West Java exhibited high complexity during the 2020--2024 period. As the center of national economic agglomeration, these three provinces experienced simultaneous changes due to the COVID-19 pandemic, followed by an uneven recovery phase across regencies/cities. These changes occurred not only in macroeconomic indicators but also propagated to dimensions of social welfare and human capital quality. To thoroughly understand the evolution of the structural relationships among these indicators, a three-way data analysis approach capable of capturing spatio-temporal interaction patterns is required.

This research applies the STATIS (\textit{Structuration des Tableaux à Trois Indices de la Statistique}) method to analyze the dynamics of 12 strategic development indicators---covering economic aspects (GRDP), welfare (Poverty, Poverty Line, Food), human development (HDI, Education, Health), labor, and equality (Gender, Gini Ratio)---across all regencies/cities in the three provinces. The analysis centers on the formation of a compromise structure. Referring to the definition by \citet{lavit94}, the compromise structure ($\mathbf{W}_c$) is constructed as a weighted sum of annual intra-structure matrices, where the weights ($\alpha_k$) are determined based on spectral decomposition to maximize inter-year similarity. This structure represents the most dominant and stable inter-regional relationship patterns during the observation period. Furthermore, this study tests the significance of the structure's stability using permutation tests and integrates the Dynamic Biplot concept to visualize the movement trajectories of regencies/cities from the crisis to the recovery period in a low-dimensional space.

The entire analysis procedure, from data pre-processing to visualization, is implemented using the Python programming language within an interactive web dashboard. This system provides a comprehensive suite of analysis features, ranging from diagnostic tests (KMO, Bartlett) and STATIS analysis to spatio-temporal cluster evaluation (Elbow Method, Dendrograms, Transition Matrices). Through interactive features such as baseline rotation, regional trajectory filtering, and automatic integration of results into LaTeX. This application serves as a sophisticated analytical and visual tool for researchers and policymakers to monitor resilience and accelerate regional development equity in the post-pandemic era.

\vspace*{0.2cm}

\noindent \bo{Keywords}:
	Interactive Visualization, K-Means Clustering, Panel Data, Regional Development, Spatial Autocorrelation, Three-way Analysis. % ------------ Jangan lupa diedit
\end{singlespace}

\newpage 